% Fragmento sugerido para documentar la prueba v0.3.
% Actualizar los valores reales después de la ejecución.

\subsection{Resultados de validación geoespacial local v0.3}

La tercera ejecución automatizada se enfocó en verificar técnicamente los productos
geoespaciales generados por el sistema después de una exportación. La prueba validó
la carpeta de salida, el archivo \texttt{polygon.json}, el video MP4, el archivo ZIP,
los GeoTIFF generados y el CSV de píxeles imputados.

\begin{table}[H]
\centering
\begin{tabular}{|p{2.4cm}|p{8.4cm}|p{2.3cm}|}
\hline
\textbf{ID} & \textbf{Validación} & \textbf{Estado} \\
\hline
GEO-00 & Directorio general de exportaciones disponible. & Pendiente \\
\hline
GEO-01 & Directorio específico de exportación disponible. & Pendiente \\
\hline
GEO-02 & Archivo \texttt{polygon.json} existente y con vértices válidos. & Pendiente \\
\hline
GEO-03 & Video MP4 existente y no vacío. & Pendiente \\
\hline
GEO-04 & ZIP GeoTIFF/CSV existente, legible y con archivos esperados. & Pendiente \\
\hline
GEO-05 & Directorio \texttt{geotiff} existente. & Pendiente \\
\hline
GEO-06 & Existencia de archivos GeoTIFF. & Pendiente \\
\hline
GEO-07 & Apertura de GeoTIFF con \texttt{rasterio} y metadatos válidos. & Pendiente \\
\hline
GEO-08 & CSV existente, legible y con columnas esperadas. & Pendiente \\
\hline
\end{tabular}
\caption{Resultados de validación geoespacial local v0.3}
\end{table}

La evidencia se almacena en \texttt{tests\_evidence/geospatial/}, incluyendo el log de
consola y el reporte JSON generado por el script de validación.
